\begin{resumo}[RESUMO]
\SingleSpacing
Este trabalho documenta o planejamento e o desenvolvimento de um sistema integrado de rastreabilidade para uma distribuidora fictícia que recebe mercadorias de fornecedores, armazena caixas e realiza entregas com frota própria. O problema central identificado foi a falta de rastreabilidade das caixas e dos veículos: a entrada, o carregamento, a saída pela portaria, a entrega e o retorno eram controlados em papel, planilhas e mensagens, o que gerava caixas extraviadas, veículos saindo com carga incompleta e pedidos sem acompanhamento. A solução é composta por uma aplicação web em ASP.NET Core MVC para gestão, logística e frota, uma API REST mínima com autenticação JWT e um aplicativo Flutter que permite ao conferente, ao operador, ao porteiro e ao motorista registrar cada etapa pela leitura de códigos de barras e QR Codes. Os três componentes compartilham uma biblioteca central em C\# com as entidades, os serviços de regra de negócio e o acesso ao SQL Server por meio do Entity Framework Core, além de \textit{triggers} e \textit{procedures} que reforçam regras críticas no próprio banco. Foram levantados dezenove requisitos funcionais e dezenove regras de negócio, todos implementados e validados por testes automatizados de ponta a ponta. O trabalho também apresenta o modelo de negócio, as ações de inclusão e acessibilidade, a proposta de implantação em nuvem com contêineres e integração contínua e o planejamento ágil do projeto.

\vspace{\onelineskip}
\noindent\textbf{Palavras-chave:} Rastreabilidade; Logística; ASP.NET Core; Flutter; SQL Server.
\end{resumo}

\begin{resumo}[ABSTRACT]
\begin{otherlanguage*}{english}
\SingleSpacing
This work documents the planning and development of an integrated traceability system for a fictional distribution company that receives goods from suppliers, stores boxes in a warehouse and delivers them with its own fleet. The core problem was the lack of traceability of boxes and vehicles: receiving, loading, gate exit, delivery and return were controlled with paper, spreadsheets and instant messages, which led to lost boxes, vehicles leaving with incomplete loads and orders without tracking. The solution consists of an ASP.NET Core MVC web application for management, logistics and fleet control, a minimal REST API with JWT authentication, and a Flutter mobile app that allows the receiving clerk, the loader, the gatekeeper and the driver to record every step by scanning barcodes and QR codes. The three components share a central C\# library containing the entities, the business rule services and SQL Server access through Entity Framework Core, in addition to triggers and stored procedures that enforce critical rules in the database itself. Nineteen functional requirements and nineteen business rules were defined, all of them implemented and validated by automated end-to-end tests. The work also presents the business model, the inclusion and accessibility actions, the cloud deployment proposal with containers and continuous integration, and the agile project plan.

\vspace{\onelineskip}
\noindent\textbf{Keywords:} Traceability; Logistics; ASP.NET Core; Flutter; SQL Server.
\end{otherlanguage*}
\end{resumo}
